\section{Main results}\label{sec:results}

\subsection{The wage response}

The central result of the paper is that the 2022 reform did not raise the real wages of
low-skilled workers relative to high-skilled workers. Column (1) of
Table~\ref{tab:main} reports the two-way fixed-effects difference-in-differences
estimate from equation~\eqref{eq:twfe}. The coefficient on the interaction of the
low-skill indicator with the post-reform indicator is $-0.009$ for the log real hourly
wage, with a standard error of $0.022$, an estimate that is economically small and
statistically indistinguishable from zero. The implied point estimate is a relative
change of less than one percent, and the ninety-five percent confidence interval rules
out relative wage gains larger than about three percent. The result for log monthly
earnings is likewise a tight null. The doubly robust estimator of
\citet{santanna_zhao_2020} in column (2) is negative and imprecise but tells the same
qualitative story of no positive wage effect.\footnote{The doubly robust estimator collapses
the data to a single pre- and post-reform period and omits department fixed effects, so it is
less comparable to the two-way fixed-effects specification and the event study; I report it as
a complementary check rather than as the preferred estimate.}

The dynamic event study confirms that this null is not an artifact of averaging over
offsetting periods. The top-left panel of Figure~\ref{fig:event} plots the coefficients
$\beta_k$ from equation~\eqref{eq:event} for the log real hourly wage. The pre-reform
coefficients are small and statistically insignificant, and a joint test of their
equality to zero does not reject, with a $p$-value of $0.60$ reported in column (3) of
Table~\ref{tab:main}. After the reform the coefficients hover around zero with no
discernible break. The wage distribution tells a consistent story:
Figure~\ref{fig:dist} shows the earnings distribution of low-skilled wage employees
before and after the reform, with the old and new minima marked. There is a visible mass
of workers near the old minimum, but the distribution does not shift appreciably toward
the new floor, and a large share of workers remain below it, reflecting the pervasive
non-compliance documented in Section~\ref{sec:data}.

\begin{table}[!tbp]\centering
\caption{Difference-in-differences estimates of the 2022 minimum-wage reform}\label{tab:main}
\begin{threeparttable}
\footnotesize
\setlength{\tabcolsep}{4pt}
\begin{tabular}{lcccc}
\toprule
Outcome & (1) TWFE DiD & (2) Doubly robust & (3) Pre-trend $p$ & (4) Elasticity \\
\midrule
Log real hourly wage & -0.021 & -0.056 & 0.001 & -0.20 \\
 & (0.012) & (0.137) & & \\
 & \multicolumn{2}{c}{\footnotesize $N=69,534$, $\bar{y}=1.744$} & & \\
Log real monthly earnings & -0.010 & -0.047 & 0.079 & -0.10 \\
 & (0.013) & (0.097) & & \\
 & \multicolumn{2}{c}{\footnotesize $N=70,691$, $\bar{y}=6.901$} & & \\
Employment & -0.025$^{***}$ & 0.020 & 0.000 & -0.35 \\
 & (0.006) & (0.037) & & \\
 & \multicolumn{2}{c}{\footnotesize $N=278,064$, $\bar{y}=0.707$} & & \\
Labour force participation & -0.008$^{*}$ & -0.015 & -- & -0.10 \\
 & (0.004) & (0.034) & & \\
 & \multicolumn{2}{c}{\footnotesize $N=278,064$, $\bar{y}=0.744$} & & \\
Formal employment & -0.046$^{***}$ & 0.006 & 0.148 & -1.99 \\
 & (0.007) & (0.040) & & \\
 & \multicolumn{2}{c}{\footnotesize $N=181,317$, $\bar{y}=0.228$} & & \\
Self-employment & 0.036$^{***}$ & 0.046 & 0.003 & 0.91 \\
 & (0.006) & (0.045) & & \\
 & \multicolumn{2}{c}{\footnotesize $N=181,317$, $\bar{y}=0.389$} & & \\
Weekly hours & 0.556$^{**}$ & -0.424 & 0.034 & 0.13 \\
 & (0.249) & (0.872) & & \\
 & \multicolumn{2}{c}{\footnotesize $N=179,111$, $\bar{y}=40.655$} & & \\
Paid below the MW & -0.006 & -0.007 & 0.128 & -0.16 \\
 & (0.007) & (0.039) & & \\
 & \multicolumn{2}{c}{\footnotesize $N=70,691$, $\bar{y}=0.352$} & & \\
\bottomrule
\end{tabular}
\begin{tablenotes}\footnotesize
\item Notes: Each row is a separate regression on ENAHO 2021Q1--2024Q4, excluding the partially treated transition quarter 2022Q2. Column (1) reports the two-way fixed-effects DiD coefficient on $\mathrm{Low}\times\mathrm{Post}$ from equation~\eqref{eq:twfe}, with department and quarter fixed effects and controls (age, age squared, sex, urban, marital status, years of education). Column (2) reports the doubly robust DiD estimator of \citet{santanna_zhao_2020} collapsing to pre/post; its standard errors are clustered by department using the estimator's influence function. Column (3) is the $p$-value of a joint test that the pre-reform event-study interactions are zero. Column (4) converts the column-(1) estimate into an elasticity with respect to the 10.2 percent minimum-wage increase (for level outcomes, relative to the baseline mean $\bar{y}$). Wage outcomes are for private-sector wage earners (employees and domestic workers, monthly-equivalent earnings); formality, hours, and self-employment are for private-sector workers. Standard errors clustered by department (25 clusters) in parentheses; all estimates are weighted by ENAHO survey weights. $^{*}p<0.1$, $^{**}p<0.05$, $^{***}p<0.01$.
\end{tablenotes}
\end{threeparttable}\end{table}


That the reform failed to lift wages at the bottom is not surprising in this setting: once
inflation is netted out the real value of the floor rose little, and most low-skilled workers
are informal or self-employed, so the statutory minimum binds for only a minority of them. The
precise null also contrasts with the ``lighthouse'' pattern that \citet{jimenez_2023} finds
for the 2016 reform. I take up these mechanisms, and the contrast with earlier reforms, in
Section~\ref{sec:discussion}.

\subsection{Employment and its composition}

If the reform did not raise pay, did it affect employment or its composition? Here the
evidence is more suggestive than conclusive, and it points to adjustment along the margin
between formal and informal work rather than to large changes in the level of employment.

The estimated effect on the probability of formal employment among the employed is
$-0.036$ (column 1 of Table~\ref{tab:main}), a decline of about three and a half
percentage points for low-skilled relative to high-skilled workers, and it is precisely
estimated under department clustering. The pre-trend test does not reject parallel trends
($p=0.44$), and the event study in Figure~\ref{fig:event} shows relative formality of the
low-skilled drifting down after the reform. Mirroring this, the probability of
self-employment rises by about three percentage points, an estimate that is statistically
significant and, as shown below, robust to randomization inference. Taken together these
two results describe a reallocation of low-skilled workers away from formal salaried jobs
and toward own-account work, the classic informalization margin emphasized by
\citet{jaramillo_lopez_2006} for Peru and by \citet{ham_2018_honduras} for Honduras. Weekly hours
in the main occupation rise modestly, consistent with a shift toward self-employment,
where hours are less constrained.

The effect on the overall employment rate is negative, at about $-2.5$ percentage points,
but it is the least robust of my estimates. As the event study makes clear, the employment
series carries a differential pre-trend concentrated in the first quarter of 2021, when the
labor market was still recovering unevenly from the pandemic across skill groups. Once that
quarter is excluded the pre-trend test no longer rejects, but the point estimate also
attenuates, and the effect does not survive my more conservative inference procedures. I
therefore treat the employment-level result as inconclusive and place little weight on it,
reserving the interpretive emphasis for the compositional shift, which is more stable.
Finally, the indicator for being paid below the minimum shows no significant change, meaning
that the reform did not measurably improve compliance at the bottom of the distribution.

\begin{figure}[H]
\centering
\includegraphics[width=0.98\textwidth]{fig5_event_study.pdf}
\caption{Event-study estimates of the 2022 reform. Each panel plots the coefficients
$\beta_k$ on the interaction of the low-skill indicator with quarter indicators from
equation~\eqref{eq:event}, relative to the omitted quarter 2022Q1 ($k=-1$). The reform is
effective at $k=0$ (2022Q2). Shaded bands are ninety-five percent confidence intervals with
standard errors clustered by department.}
\label{fig:event}
\end{figure}

\begin{figure}[H]
\centering
\includegraphics[width=0.72\textwidth]{fig4_wage_distribution.pdf}
\caption{Earnings distribution of low-skilled wage employees, before and after the reform.
Weighted kernel densities of nominal monthly earnings, with the old (930) and new (1,025)
minimum wage marked by dashed lines.}
\label{fig:dist}
\end{figure}
